\documentclass[11pt]{article}

\usepackage[T1]{fontenc}
\usepackage[utf8]{inputenc}
\usepackage{amsmath,amssymb,amsthm}
\usepackage{booktabs}
\usepackage{geometry}
\usepackage{hyperref}
\usepackage{microtype}
\usepackage{xcolor}

\geometry{margin=1in}
\hypersetup{
  colorlinks=true,
  linkcolor=blue!55!black,
  citecolor=blue!55!black,
  urlcolor=blue!55!black
}

\newtheorem{theorem}{Theorem}[section]
\newtheorem{proposition}[theorem]{Proposition}
\newtheorem{corollary}[theorem]{Corollary}
\theoremstyle{definition}
\newtheorem{definition}[theorem]{Definition}
\newtheorem{conjecture}[theorem]{Conjecture}
\theoremstyle{remark}
\newtheorem{remark}[theorem]{Remark}

\newcommand{\F}{\mathbb F}
\newcommand{\supp}{\operatorname{supp}}
\newcommand{\symdiff}{\mathbin{\triangle}}
\newcommand{\FiveCDC}{\textup{FiveCDC}}

\title{Minimum Extendable Fano Projections and the\\
Five-Cycle Double Cover Conjecture}
\author{Atharva Vaidya\\
\small AI-assisted research note for independent human verification}
\date{July 29, 2026}

\begin{document}
\maketitle

\begin{center}
\fcolorbox{red!65!black}{red!4}{
\parbox{0.91\textwidth}{
\textbf{Resolution status.}
The five-cycle double cover conjecture remains unresolved.
This note proves an exchange theorem, settles the minimum-support cases
through size seven, and gives exact finite censuses.  It proves neither the
conjecture nor its negation.
}}
\end{center}

\begin{abstract}
Hu\v{s}ek and \v{S}\'amal recently expressed the five-cycle double cover
conjecture as the problem of selecting a nowhere-zero
\(\F_2^3\)-flow with a component-parity property.  We study the support
\(h\) of one binary coordinate of such a flow.  Call \(h\)
\emph{extendable} if it occurs as a coordinate of some nowhere-zero
\(\F_2^3\)-flow.

We prove a simultaneous exchange theorem for a minimum-cardinality
extendable projection.  Write the flow as \(f=(h,s)\), with
\(s:E\to\F_2^2\), and partition \(h\) into the four affine value classes
\(M_c=\{e\in h:s(e)=c\}\).  For every \(c\), the projection \(h\) is a
minimum-cardinality binary cycle containing \(M_c\).  Equivalently, every
binary cycle \(C\) avoiding \(M_c\) satisfies
\[
                         |C\cap h|\le |C-h|.
\]
Thus every circuit component of \(h\) contains all four affine colours,
and every support-shortening cycle meets all four classes.  The proof is
the addition of the Fano value \((1,c)\) along \(C\).

We first prove, without computation, that every globally minimum
extendable projection of size at most five is cleanable.  If such an
extension were dirty, its projection would be a single four- or five-cycle.
Classifying the possible component shores and applying independent
\(\operatorname{GL}(2,2)\) transformations on the two complementary
components constructs a nowhere-zero \(\F_2^2\)-flow, contradicting
non-Tait-colourability.  A second exact theorem settles sizes six and
seven.  An
exhaustive affine/dihedral boundary quotient either gives the same Tait
repair or leaves only states with two complement components, all cleaned
by a proved one-switch span theorem.  Thus a
counterexample to the proposed minimum-projection selection principle must
have minimum size at least eight.

We perform an exact canonical census of all \(41{,}301\) connected simple
cubic graphs of order \(18\).  Among the \(39{,}866\) bridgeless graphs,
exactly \(179\) are not Tait-colourable, and every minimum-cardinality
extendable projection on all \(179\) is cleanable.  A second exact replay
uses frozen canonically generated sources: \(14{,}009\) cyclically
4-edge-connected non-Tait records through order \(28\), and a separate
\(12{,}892\)-record hard order-\(22\) source.  All \(147{,}539\) minimum
extendable projections in those files are cleanable.  Completeness of the
source populations is inherited from their upstream generation packages;
the result on the literal frozen files is checked directly.
An additional exact SAT scan covers seven retained order-\(34\) and
\(7{,}654\) retained order-\(40\) strong-snark rows.  All \(433{,}730\)
minimum projections in those files have size ten and are cleanable.

The universal statement suggested by this computation is recorded as an
open conjecture.  The missing step is to derive a colour-avoiding strict
exchange from a rainbow-odd component, possibly after support-neutral
recolouring.  We make no claim of a FiveCDC resolution or of priority for
the flow formulation.
\end{abstract}

\tableofcontents

\section{Status, conventions, and contribution boundary}

A binary cycle of a graph is an Eulerian edge set; it may be empty or
disconnected.  A \(k\)-cycle double cover is an indexed family of \(k\)
binary cycles in which every edge occurs exactly twice.  Empty members may
be appended, so ``at most five'' and ``five, allowing empty members'' are
equivalent.  The five-cycle double cover conjecture, attributed to Celmins
and Preissmann, asserts that every bridgeless graph has such a cover with
at most five members~\cite{Celmins1985,Preissmann1981}.

We work with the cubic flow formulation in
\cite[Conjecture~3.19]{HusekSamal2026}.  No general-to-cubic reduction is
proved in this note.  All theoretical statements below apply to finite
loopless cubic multigraphs unless ``simple'' is stated.  Parallel edges
are distinct edge objects.  A loop would contribute twice to vertex
degree; loops are excluded here to keep the cut notation literal.  The
finite census is restricted to simple graphs.

Hu\v{s}ek and \v{S}\'amal prove that the cubic FiveCDC problem is equivalent
to selecting a nowhere-zero \(\F_2^3\)-flow with an additional parity
condition~\cite[Theorem~3.16 and Conjecture~3.19]{HusekSamal2026}.  Their
criterion is related to the prescribed-cycle formulation of
Hoffmann-Ostenhof~\cite{HoffmannOstenhof2013}.  We claim no novelty or
priority for either formulation.

The contribution of this note is narrower:
\begin{enumerate}
  \item a minimum-support exchange theorem, with a complete proof;
  \item a human proof through size five and an exact extension through
        size seven;
  \item four simultaneous shortest-join consequences;
  \item a reproducible canonical order-\(18\) census and an expanded
        exact replay on frozen sources through order \(28\); and
  \item a precise conjectural step that would turn this route into a proof.
\end{enumerate}
The last item is open.

\section{Extendable and clean Fano projections}

Let \(G=(V,E)\) be a loopless cubic graph and put \(K=\F_2^2\).
Write a three-bit flow as
\[
                 f=(h,s):E\longrightarrow \F_2\times K,       \tag{1}
\]
where \(h:E\to\F_2\) and \(s:E\to K\) are flows.  We identify a binary
flow with its support.

\begin{definition}
A binary cycle \(h\), possibly zero, is \emph{extendable} if there is a flow
\(s:E\to K\) for which \(f=(h,s)\) is nowhere zero.
\end{definition}

Writing \(s=(p,q)\), extendability has the following elementary form:
\[
                         E-h\subseteq p\cup q,                 \tag{2}
\]
where \(p\) and \(q\) are binary cycles.  Indeed, an edge of \(E-h\) has
first coordinate zero and therefore needs a nonzero low coordinate.

Fix an extension \(f=(h,s)\).  For \(c\in K\), put
\[
                         M_c=\{e\in h:s(e)=c\}.                \tag{3}
\]

\begin{proposition}\label{prop:matching}
Every \(M_c\) is a matching, and the four sets \(M_c\) partition \(h\).
\end{proposition}

\begin{proof}
The partition statement is immediate.  If two edges incident with a cubic
vertex both had flow value \((1,c)\), their sum would be zero and flow
conservation would force the third incident edge to have value zero,
contrary to nowhere-zeroness.
\end{proof}

Let \(W\) be a component of the spanning subgraph \((V,E-h)\).  Every edge
of \(\delta(W)\) belongs to \(h\).  If
\[
                    n_c(W)=|\delta(W)\cap M_c|\pmod 2,
\]
then cut conservation for the first coordinate and the two low coordinates
gives
\[
              \sum_{c\in K}n_c(W)=0,\qquad
              \sum_{c\in K}n_c(W)c=0.                         \tag{4}
\]
The \(3\times4\) binary matrix in (4) has one-dimensional kernel, generated
by the all-ones vector.  Hence the four parities \(n_c(W)\) are equal.

\begin{definition}
An extension is \emph{clean} if \(n_c(W)=0\) for every component \(W\)
of \((V,E-h)\), for one, and hence every, \(c\in K\).
A projection \(h\) is \emph{cleanable} if it has a clean extension.
A component with all four parities equal to one is \emph{rainbow-odd}.
\end{definition}

If \(s=(p,q)\), the class \(M_{11}\) on \(\delta(W)\) is
\(\delta(W)\cap p\cap q\).  Thus cleanability is exactly the existence of
binary cycles \(p,q\) satisfying (2) and
\[
             |\delta(W)\cap p\cap q|\equiv0\pmod2             \tag{5}
\]
for every component \(W\) of \((V,E-h)\).  By
\cite[Theorem~3.16]{HusekSamal2026}, existence of such an \(h\) is
equivalent to the cubic FiveCDC assertion.  In the Tait-colourable case one
may take \(h=\varnothing\) and a nowhere-zero \(K\)-flow.

\section{The minimum-projection exchange theorem}

Assume now that \(G\) is not Tait-colourable, so an extendable projection
cannot be empty.  Choose \(h\) of minimum cardinality among all extendable
projections, and fix any extension \(f=(h,s)\).

\begin{theorem}[minimum-projection exchange]\label{thm:exchange}
For every \(c\in K\) and every binary cycle \(C\) with
\(C\cap M_c=\varnothing\),
\[
                         |C\cap h|\le |C-h|.                    \tag{6}
\]
Equivalently, \(h\) is a minimum-cardinality binary cycle containing
\(M_c\), simultaneously for all four \(c\in K\).
\end{theorem}

\begin{proof}
Fix \(c\in K\), put \(t=(1,c)\), and add \(t\) to the flow value on every
edge of \(C\):
\[
 f'(e)=
 \begin{cases}
   f(e)+t,&e\in C,\\
   f(e),&e\notin C.
 \end{cases}                                                   \tag{7}
\]
Since \(C\) is a binary cycle, the added function is a flow.  The only
edges that could become zero in (7) are those whose old value was \(t\),
and these are exactly the edges of \(M_c\).  The hypothesis
\(C\cap M_c=\varnothing\) therefore makes \(f'\) nowhere zero.

The first coordinate of \(f'\) is \(h\symdiff C\).  This is another
extendable projection, so minimality gives
\[
 |h|\le |h\symdiff C|
     =|h|+|C-h|-|C\cap h|.
\]
This is (6).

For the equivalent statement, rebase the low coordinates by replacing
\(s\) with the flow \(s+ch\).  Its exact zero set is \(M_c\): outside
\(h\), the low value was already nonzero, while on \(h\) it vanishes
exactly when \(s(e)=c\).

Let \(h'\) be any binary cycle containing \(M_c\).  The flow
\((h',s+ch)\) is nowhere zero.  On an edge outside \(h'\), its low value is
nonzero because every zero lies in \(M_c\subseteq h'\); on an edge of
\(h'\), its first coordinate is one.  Thus \(h'\) is extendable, and
\(|h|\le |h'|\).

Conversely, if \(C\cap M_c=\varnothing\), then
\(h\symdiff C\) contains \(M_c\).  Applying the minimum-containing
property to \(h\symdiff C\) recovers (6).
\end{proof}

\subsection{Shortest joins and four-colour consequences}

Let \(T_c\) be the set of endpoints of \(M_c\).  Removing \(M_c\) from a
binary cycle containing \(M_c\) produces a \(T_c\)-join in \(G-M_c\), and
adjoining \(M_c\) reverses the correspondence.  Therefore
Theorem~\ref{thm:exchange} has the following form.

\begin{corollary}\label{cor:tjoin}
For every \(c\in K\), the edge set
\[
                              J_c=h-M_c
\]
is a shortest \(T_c\)-join in \(G-M_c\).
\end{corollary}

This is not the already-refuted principle that a cardinality-minimum
exact-zero matching must extend to a FiveCDC certificate.  Here the
objective is the cardinality of the whole binary cycle
\(M_c\cup J_c\), and the same projection is optimal after all four affine
rebases.

\begin{corollary}[flow-resistance bound]\label{cor:resistance}
Let \(r_f(G)\) be the minimum number of zero edges in an
\(\F_2^2\)-flow on \(G\).  Every extendable projection satisfies
\[
                              |h|\ge4r_f(G).                    \tag{8}
\]
\end{corollary}

\begin{proof}
For each \(c\), the rebased low flow \(s+ch\) has exact zero set \(M_c\).
Hence \(|M_c|\ge r_f(G)\).  The four \(M_c\)'s partition \(h\).
\end{proof}

\begin{corollary}[four-colour circuits]\label{cor:fourcolour}
Every circuit component of \(h\) contains at least one edge of each of the
four sets \(M_c\).
\end{corollary}

\begin{proof}
If a circuit component \(D\) omitted \(M_c\), then \(D\) would be a binary
cycle avoiding \(M_c\), but
\(|D\cap h|=|D|>|D-h|=0\), contrary to (6).
\end{proof}

\begin{theorem}[minimum projections through size five]
\label{thm:sizefive}
Let \(G\) be a connected bridgeless loopless cubic graph; parallel edges
are allowed.  If \(G\) is not Tait-colourable and its minimum extendable
projection has size at most five, then every minimum extendable projection
is cleanable.  If \(G\) is Tait-colourable, the zero projection is a clean
minimum.
\end{theorem}

\begin{proof}
The Tait case follows by taking a nowhere-zero \(K\)-flow as the low
coordinate: its zero affine class is empty, so the equal cut parities in
(4) are all even.  Assume henceforth that \(G\) is non-Tait, choose a
minimum extendable \(h\), and fix an arbitrary extension \(f=(h,s)\).

By Corollary~\ref{cor:fourcolour}, every circuit component of \(h\) meets
all four \(M_c\).  Since \(G\) is loopless, \(|h|\le5\) therefore forces
\(h\) to be one ordinary circuit \(D\) of length \(n=4\) or \(5\); a
parallel two-circuit cannot meet four classes.  Write
\[
 v_0,e_0,v_1,e_1,\ldots,v_{n-1},e_{n-1},v_0,
 \qquad e_i=v_iv_{i+1},
\]
cyclically, and put \(c_i=s(e_i)\).  Adjacent \(c_i\)'s are distinct:
otherwise the third edge at their common vertex would have full flow
value zero.

Let \(K_0=G-h\).  Every component of \(K_0\) contains a vertex of \(D\),
by connectedness.  Suppose the extension is dirty.  For a dirty component
\(W\), membership along \(D\) toggles oddly on each of the four affine
colour classes.

If \(n=4\), the four colours occur once each, so membership toggles on
every \(e_i\).  The components of \(K_0\) are exactly the two alternating
vertex pairs.  Indeed, separating either pair further would create a
component whose cut sees exactly two affine colours oddly, contrary to
(4).

If \(n=5\), rotate and rename the colours so their cyclic word is
\(A,B,A,C,D\).  Put
\[
 z_i=1[v_i\in W],\qquad y_i=z_i+z_{i+1}.
\]
Dirtiness says
\[
 y_1=y_3=y_4=1,\qquad y_0+y_2=1.
\]
Up to complement, the solutions are \(W=\{v_2,v_4\}\) and
\(W=\{v_1,v_4\}\).  Summing the common odd cut parity over all
\(K_0\)-components shows that the number of dirty components is even.
The only component-compatible shore disjoint from either displayed dirty
shore is its complement: exhaustive substitution in the preceding four
equations also gives the clean pair \(\{v_1,v_2\}\) and its complement,
both of which intersect either dirty partition.  Hence \(K_0\) again has
exactly two components.

Let \(k_i\) be the unique \(K_0\)-edge incident with \(v_i\), and put
\[
 d_i=s(k_i)=c_{i-1}+c_i\ne0.
\]
On each of the two \(K_0\)-components, apply an independently chosen
element of \(\operatorname{GL}(2,2)\) to every low flow value.  This
preserves internal conservation and nonzeroness.  We now choose nonzero
values \(r_i\) on \(e_i\) so that
\[
                         r_{i-1}+r_i=g_i,                 \tag{9}
\]
where \(g_i\) is the transformed value on \(k_i\).

For \(n=4\), the boundary values on each alternating pair agree:
\(d_0=d_2=x\) and \(d_1=d_3=y\).  Map \(x\) and \(y\), independently, to
the same nonzero \(\alpha\).  If
\(\{\alpha,\beta,\gamma\}=K-\{0\}\), assigning \(\beta,\gamma\)
alternately to the edges of \(D\) satisfies (9).

For \(n=5\), translate the affine colour names so \(A=0\), write
\(B=x,C=y,D=x+y\), and normalize \(x=1,y=2,x+y=3\).  Thus
\[
                         d=(3,1,1,2,1).
\]
The complete repairs for the two component partitions are
\[
\begin{array}{c|c|c}
K_0\text{-components on }v_i& (g_0,\ldots,g_4)&(r_0,\ldots,r_4)\\
\hline
\{0,1,3\}\mid\{2,4\}&(3,1,2,2,2)&(2,3,1,3,1)\\
\{0,2,3\}\mid\{1,4\}&(3,3,1,2,3)&(1,2,3,1,2).
\end{array}
\]
In each row the three-terminal component uses the identity.  On the
two-terminal component, both original boundary values are \(1\), which an
invertible linear map sends respectively to \(2\) or \(3\).  Direct
cyclic substitution verifies (9).

We have constructed a nowhere-zero \(K\)-flow on every edge of \(G\).
At a cubic vertex its three incident values are therefore the three
distinct nonzero elements of \(K\), giving a Tait colouring.  This
contradicts the hypothesis.  Thus the arbitrary extension of \(h\) was
clean, and in particular \(h\) is cleanable.
\end{proof}

\begin{remark}
The finite case split and the two repair rows are independently replayed
by \path{verify.py} in the bundled size-five theorem package.
\end{remark}

\subsection{A line-switch lemma and the size-six/seven boundary}

We need one support-preserving fact.  Identify the seven nonzero full flow
values with \(1,\ldots,7\) under bitwise addition, so that
\[
 L=\{1,2,3\}=\{(0,t):t\in K-\{0\}\},\qquad
 A=\{4,5,6,7\}=\{(1,c):c\in K\}.
\]
Let \({\cal K}\) be the components of \(G-h\), and let
\(r\in\F_2^{\cal K}\) indicate the rainbow-odd components.  For
\(t\in L\), put \(N_t=f^{-1}(t)\) and \(P_t=\{4,4+t\}\).  If \(C\) is a
binary cycle of \(G-N_t\), define
\[
 \tau_t(C)_W
   =|C\cap\delta(W)\cap f^{-1}(P_t)|\pmod2
   \qquad(W\in{\cal K}).                                \tag{10}
\]
Adding \(t\) along \(C\) is a valid line-preserving switch: avoiding
\(N_t\) prevents a zero value, and (10) is its change to the defect
vector.

\begin{proposition}[combined-line span]\label{prop:span}
\[
 r\in\sum_{t\in L}\operatorname{im}\tau_t.              \tag{11}
\]
Consequently, if \(G-h\) has at most two components, then \(h\) is
cleanable by at most one valid line-preserving switch.
\end{proposition}

\begin{proof}
Suppose (11) fails.  There is a vector
\(y\in\F_2^{\cal K}\) with
\[
 y\cdot r=1,\qquad y\cdot\tau_t(C)=0
 \quad(t\in L,\ C\text{ a cycle of }G-N_t).
\]
Let \(U\) be the union of the components selected by \(y\), and let
\(u:V(G)\to\F_2\) indicate \(U\).  Cycle--cut orthogonality says that
\[
 R_t=\delta(U)\cap f^{-1}(P_t)
\]
is a cut in \(G-N_t\).  Choose a vertex indicator \(x_t\) for that cut,
and write \(q=(u,x_1,x_2,x_3)\).

Define
\[
 {\cal A}(q)=
 \bigl(u(x_2+x_3),\ u(x_1+x_3),\
 x_1x_2+x_1x_3+x_2x_3\bigr)\in\F_2^3.
\]
Across an edge, the only possible nonzero differences of
\((u;x_1x_2x_3)\), indexed by its flow value, are
\[
\begin{array}{c|ccccccc}
f(e)&1&2&3&4&5&6&7\\
\hline
\Delta q&(0;100)&(0;010)&(0;001)&(1;111)&
(1;100)&(1;010)&(1;001).
\end{array}
\]
The zero difference is also allowed in every column.  Direct substitution
gives the pointwise identity
\[
 \bigl({\cal A}(q_v)+{\cal A}(q_w)\bigr)\cdot f(vw)
 =
 \begin{cases}
 1,&f(vw)=4\text{ and }u(v)\ne u(w),\\
 0,&\text{otherwise}.
 \end{cases}                                           \tag{12}
\]
Summing (12) over all edges and regrouping at vertices makes its left side
\[
 \sum_v{\cal A}(q_v)\cdot\left(\sum_{e\ni v}f(e)\right)=0
\]
by flow conservation.  The right side is
\(|\delta(U)\cap f^{-1}(4)|=y\cdot r=1\), a contradiction.  This proves
(11).

Every vector in every image of \(\tau_t\) has even Hamming weight, since
an affine edge between two \(G-h\) components is counted at both ends.
For one component the even-weight space is zero.  For two components it
is one-dimensional; if \(r\ne0\), (11) forces one of the three image
spaces to contain its unique nonzero vector.  The corresponding single
valid switch kills \(r\).  A single switch has no mixed-switch quadratic
correction, completing the proof.
\end{proof}

\begin{theorem}[minimum projections through size seven]
\label{thm:sizeseven}
Under the hypotheses of Theorem~\ref{thm:sizefive}, every minimum
extendable projection of size at most seven is cleanable.
\end{theorem}

\begin{proof}
Only sizes six and seven remain.  The four-colour circuit corollary forces
\(h\) to be one circuit
\[
 v_0,e_0,v_1,e_1,\ldots,v_{n-1},e_{n-1},v_0,
 \qquad n\in\{6,7\}.
\]
Put \(c_i=s(e_i)\), and encode the component of \(G-h\) containing \(v_i\)
by \(\pi_i\).  The \(c_i\)'s form a proper cyclic word using all four
affine colours.  Conservation on every component \(W\) gives
\[
 \sum_{\pi_i=W}(c_{i-1}+c_i)=0,
\]
and its four affine cut parities are either all zero or all one.

The exact finite boundary classifier enumerates every such word and every
set partition \(\pi\).  For each
state it also enumerates one independent element of
\(\operatorname{GL}(2,2)\) on the low flow values of every component and
solves the \(n\) cycle equations.  Its complete counts are
\[
\begin{array}{c|r|r|r}
n&\text{valid dirty pairs}&\text{failed linear repairs}&
\text{failures with more than two components}\\
\hline
6&2712&432&0\\
7&26880&5376&0.
\end{array}
\]
Every dirty state either produces a
nowhere-zero \(K\)-flow on all of \(G\), contradicting the non-Tait
hypothesis, or has exactly two components of \(G-h\).  At size six the
failed states may additionally be quotiented by the affine group on the
four colours, the dihedral group of the cycle, and component relabelling;
they give the three representatives
\[
\begin{array}{c|c}
(c_0,\ldots,c_5)&(\pi_0,\ldots,\pi_5)\\
\hline
010123&001001\\
010232&001010\\
012013&000101.
\end{array}
\]
with 144 labeled states per orbit.  At size seven there is no residual
normal form after the two-component theorem.  The dependency-free
enumerator regenerates all words, partitions, conservation tests, dirty
states, and linear repairs rather than trusting the stored counts.

Proposition~\ref{prop:span} makes every failed state clean by one valid
switch.  Hence no dirty size-six or size-seven global minimum is
uncleanable.
\end{proof}

\begin{remark}
The exact classifications, proof of semantics, size-six representatives,
and full standalone span proof are bundled in the cumulative through-seven
package.  This is a finite boundary theorem, not a search over sampled
graphs.
\end{remark}

\begin{corollary}[rainbow shortcut]\label{cor:shortcut}
Every binary cycle \(C\) satisfying \(|C\cap h|>|C-h|\) meets all four
sets \(M_c\).
\end{corollary}

\begin{proof}
If \(C\) omitted one class \(M_c\), it would contradict (6).
\end{proof}

\section{The exact missing step}

Theorem~\ref{thm:exchange} suggests the following selection principle.

\begin{conjecture}[minimum extendable projection]\label{conj:minimum}
Every bridgeless cubic graph has a minimum-cardinality extendable
projection that is cleanable.
\end{conjecture}

The stronger statement that every minimum projection is cleanable also
survives the finite census in Section~\ref{sec:census}, but it is not needed.
Conjecture~\ref{conj:minimum}, together with the Hu\v{s}ek--\v{S}\'amal
equivalence, would imply the cubic FiveCDC conjecture.

Suppose a minimum projection \(h\) is not cleanable.  Every extension then
has a rainbow-odd component \(W\) of \(G-h\), so all four \(M_c\)'s cross
\(\delta(W)\) oddly.  A direct proof of
Conjecture~\ref{conj:minimum} would need to derive either
\begin{enumerate}
  \item a binary cycle \(C\) avoiding some \(M_c\) and satisfying
        \(|C\cap h|>|C-h|\), contradicting
        Theorem~\ref{thm:exchange}; or
  \item a sequence of support-neutral low-coordinate recolourings followed
        by such a strict exchange.
\end{enumerate}

The second possibility cannot currently be discarded.  Small exact
examples have four disjoint colour classes meeting every immediately
component-reducing cycle, while a support-preserving recolouring escapes
the trap.  Other examples show that a fixed nowhere-zero Fano flow can
have all seven of its functional projections dirty.  These are
proof-strategy obstructions, not FiveCDC counterexamples.  They rule out
an argument that fixes an arbitrary initial flow or assumes one elementary
switch always suffices.

The unresolved human proof obligation is precise: from a rainbow-odd
component of an \emph{uncleanable globally minimum-support projection},
derive a negative colour-avoiding cycle after allowable neutral
recolouring, using the four simultaneous shortest-\(T_c\)-join
inequalities of Corollary~\ref{cor:tjoin}.  No proof of this statement is
given here.

\section{Exact finite census}\label{sec:census}

\subsection{Canonical order-\(18\) method}

The accompanying checker uses canonical generation rather than sampling.
For order \(18\), the command
\[
\texttt{geng -cq -d3 -D3 18}
\]
generates every connected simple cubic graph once.  The checker then:
\begin{enumerate}
  \item tests bridgelessness by deleting every edge and checking
        connectivity;
  \item tests Tait colourability by exhaustive backtracking;
  \item enumerates the entire binary cycle space;
  \item processes nonzero projections \(h\) by increasing cardinality;
  \item quotients pairs of cycles \((p,q)\) by their exact semantic state
        \((p\cup q,p\cap q)\);
  \item tests extendability by \(E-h\subseteq p\cup q\); and
  \item computes the components of \(G-h\) and tests (5) on every
        component.
\end{enumerate}

For a connected cubic graph of order \(18\), the binary cycle-space
dimension is
\[
                  |E|-|V|+1=27-18+1=10.
\]
Thus all \(2^{10}=1024\) binary cycles are enumerated.  Equal
union/intersection pairs have exactly the same coverage and component
parity behaviour, so the semantic quotient loses no information.
No snark or cyclic-connectivity reduction is used.

\subsection{Results and reproducibility}

\begin{table}[h]
\centering
\begin{tabular}{lr}
\toprule
class & count\\
\midrule
connected simple cubic graphs & \(41{,}301\)\\
bridgeless graphs & \(39{,}866\)\\
bridgeless non-Tait graphs & \(179\)\\
\bottomrule
\end{tabular}
\caption{Complete order-\(18\) population.}
\end{table}

Every minimum-cardinality extendable projection of every one of the
\(179\) non-Tait graphs is cleanable.  The minimum-size histogram is
\[
             5:166,\qquad 6:11,\qquad 7:1,\qquad 8:1.
\]
The checker hashes, in canonical graph order, the graph6 string, minimum
size, number of minimum extendable projections, and number of cleanable
minimum projections.  The resulting SHA-256 is
\begin{center}
\texttt{82b01cc2f01d4f50f7745f4bdb1b3dc46ade798d452a9d4f11d4a91d738a7834}.
\end{center}

From the project root, run:
\begin{verbatim}
python3 scratch/check_husek_samal_minimum_projection_frontier.py \
  --order 18
\end{verbatim}
The checker requires Python, the sibling cycle-space module
\path{scratch/search_fano_all_bad_projection_subspaces.py}, and
nauty \texttt{geng} on \texttt{PATH}.  It has no SAT solver, network, or
nonstandard Python-package dependency.  The relevant source hashes are
\begin{center}
\small
frontier checker:\\
\texttt{37e972782863c7a9506f6d339159ddaca3865363107704f86365c1a9b77a30c3}\\[2pt]
cycle-space module:\\
\texttt{7e5ea2599e7533f041ea88de5c6f5c7a3bf542cedc33d78728439d43ce8861e8}
\end{center}

This is a finite theorem.  It is evidence for
Conjecture~\ref{conj:minimum}, not a proof of it.

\subsection{Expanded frozen-source replay}

A separate C++ classifier implements the same extendability and
cleanability semantics by Gaussian elimination over \(\F_2\).  It replays
two previously frozen graph6 populations:
\begin{enumerate}
  \item \(14{,}009\) cyclically 4-edge-connected non-Tait records of even
        orders \(10\) through \(28\); and
  \item \(12{,}892\) bridgeless non-Tait rows from a broader hard
        order-\(22\) source.
\end{enumerate}
The exact results are:
\begin{table}[h]
\centering
\begin{tabular}{lrrr}
\toprule
frozen source & graphs & minimum projections & uncleanable\\
\midrule
cyclically 4, orders \(10\)--\(28\) & \(14{,}009\) & \(92{,}754\) & \(0\)\\
hard order \(22\) & \(12{,}892\) & \(54{,}785\) & \(0\)\\
\midrule
total & \(26{,}901\) & \(147{,}539\) & \(0\)\\
\bottomrule
\end{tabular}
\caption{Exact minimum-extendable-projection replay on frozen sources.}
\end{table}

For the cyclically 4 source the minimum-size profile is
\[
                       5:14007,\qquad 6:1,\qquad 9:1.
\]
For the hard order-\(22\) source it is
\[
 5:11776,\quad 6:1025,\quad 7:55,\quad 8:20,\quad 9:6,\quad 10:10.
\]
These counts refer to graphs; the third column of the table counts all
minimum projections on those graphs.

From the project root, run:
\begin{verbatim}
cd scratch/minimum-projection-census-through28-20260729
python3 verify_summary.py
\end{verbatim}
The replay compiles the classifier, verifies every input digest and every
emitted graph6 row, and checks all aggregates above.  Its checksum ledger is
the \texttt{SHA256SUMS} file in that package.
The unconditional claim here concerns the literal frozen files.  Their
identification with complete graph classes uses the documented upstream
canonical-generation and filtering packages; this replay does not
independently regenerate those populations.

\subsection{Retained strong snarks of orders \(34\) and \(40\)}

A separate exact scan covers the literal retained strong-snark files at
orders \(34\) and \(40\):
\begin{table}[h]
\centering
\begin{tabular}{lrrr}
\toprule
order & graph6 rows & minimum projections & uncleanable\\
\midrule
34 & \(7\) & \(371\) & \(0\)\\
40 & \(7{,}654\) & \(433{,}359\) & \(0\)\\
\midrule
total & \(7{,}661\) & \(433{,}730\) & \(0\)\\
\bottomrule
\end{tabular}
\caption{Exact minimum-projection scan on the frozen strong-snark rows.}
\end{table}
Every row has minimum extendable-projection size ten.  An ordinary-CNF
solver encodes the two low binary cycles, coverage on \(E-h\), product
variables for \(p_e\wedge q_e\), and an XOR chain imposing even product
parity on every component cut.  Every even projection through weight
twelve is enumerated as a vertex-disjoint union of circuits, which is
exhaustive for a cubic graph.  All rows reach their first extendable level
at weight ten.

The order-\(34\) output and the first five order-\(40\) rows were reproduced
independently by the linear-algebra scanner used in the earlier census.
The full order-\(40\) claim uses the exact SAT implementation.  The
concatenated result digest is
\begin{center}
\texttt{becb7e4b648c000ec874509bccdf389df24f6f0d82bce2da49ba10d5b10d38e3}.
\end{center}
Inputs, frozen output, both scanner implementations, and the replay script
are under \path{scratch/}, in the subdirectory
\path{minimum-projection-known-strong-snarks-20260729/}.
As before, the unconditional statement concerns the literal bundled rows;
completeness of the named graph populations is inherited from upstream
provenance.

\subsection{A certified 130-vertex parameter-separation test}

The retained 130-vertex exchange graph gives a substantially larger exact
stress test.  A three-coordinate CNF uses binary-cycle variables
\((h_e,p_e,q_e)\), the three parity equations at every vertex, the
nowhere-zero clauses \(h_e\vee p_e\vee q_e\), and a sequential cardinality
counter.  Its certified results are
\[
 \min |h|=42,\qquad
 \#\{h:|h|=42,\ h\text{ extendable}\}=11264,
\]
and every one of those 11,264 supports is cleanable.  The size-\(\le41\)
CNF has 7,760 variables and 16,218 clauses.  A second CNF blocks every
listed size-42 support and proves enumeration completeness.  Both UNSAT
LRATs are accepted by C \texttt{lrat-check} and verified CakeML
\texttt{cake\_lpr}; every positive clean lift is also checked directly.

On the same graph the different parameter
\[
 \rho_3(G)=
 \min_{f,\ a\ne0}|f^{-1}(a)|
\]
equals five, while no size-five value class packs two boundary joins.
Thus minimizing one nonzero Fano value class and minimizing an entire
coordinate projection are not interchangeable.  The former selection
principle fails on this graph; the latter passes all 11,264 minimum cases.
To replay the package, run:
\begin{verbatim}
cd search/minimum-projection-n130-20260729
python3 verify.py
\end{verbatim}

\section{Discussion}

The exchange theorem converts global support minimality into four ordinary
geodesic conditions: one shortest join for each affine value class.
Theorems~\ref{thm:sizefive} and~\ref{thm:sizeseven} close the first four
possible nonzero support sizes: a counterexample to
Conjecture~\ref{conj:minimum} must have minimum projection size at least
eight.
The useful feature is simultaneity.  A cycle that would shorten the
projection is allowed to be blocked by one colour, but
Corollary~\ref{cor:shortcut} forces it to be blocked by all four.
Likewise, disconnected projections are severely constrained because every
circuit component must contain every colour.

The four-colour simultaneity cannot be replaced by a theorem about one
matching.  In the Petersen graph with edge set
\[
\{03,06,09,14,16,18,25,26,27,37,38,47,49,58,59\},
\]
the matching \(M=\{06,14,25\}\) is contained in exactly two
minimum-cardinality binary cycles:
\[
\{06,09,14,18,25,26,49,58\},\qquad
\{06,09,14,16,25,27,47,59\}.
\]
Both have size eight, and in both cases the components of the complement
have \(M\)-endpoint parities \(1,0,1\).  Exhausting all \(2^{15}\) edge
subsets proves minimality.  The complete finite replay is in
\path{scratch/petersen-minimum-tjoin-countermodel-20260729/}.  This
obstruction does not realize four synchronized affine classes.

What is not yet understood is how the common rainbow-odd cut interacts with
the four shortest-join duals.  A publishable resolution would need either a
cut/cycle contradiction, a neutral-reconfiguration theorem, or an exact
counterexample to Conjecture~\ref{conj:minimum}.  Until one of those is
supplied, the correct status is a frontier result.

\section*{AI-use and verification disclosure}

This research route, proof draft, checker, computation, and exposition were
developed by OpenAI Codex agents under Atharva Vaidya's direction.  The
minimum-projection exchange proof, the size-at-most-five proof, and the
span proof used at sizes six and seven are displayed in full so that they
can be checked without trusting an AI system.  Those boundary
classifications are also regenerated by a dependency-free exact checker.
The order-\(18\) result is
reproducible from canonical generation and exhaustive cycle-space
enumeration.  The expanded census exactly replays frozen sources but
inherits their population provenance.  Both remain computational results
and should be independently rerun.

Hu\v{s}ek and \v{S}\'amal's flow criterion and
Hoffmann-Ostenhof's matching/join formulation are prior work.  This draft
makes no claim of a FiveCDC resolution and no priority claim for the
minimum-projection viewpoint pending expert literature review.  Before
public submission, a human graph theorist should check the proof line by
line, rerun the census in a clean environment, audit the bibliography, and
decide authorship and venue policy under the applicable AI-disclosure
rules.

\small
\bibliographystyle{plain}
\bibliography{references}

\end{document}
